\chapter{Justificaci\'on}

El desarrollo del presente proyecto se justifica desde una perspectiva t\'ecnica, acad\'emica y funcional.

Desde el punto de vista t\'ecnico, la disponibilidad de datasets p\'ublicos anotados y de modelos de segmentaci\'on m\'edica permite abordar problemas que anteriormente requer\'ian acceso exclusivo a datos privados o infraestructura cl\'inica compleja. En particular, el dataset SPIDER ofrece una base adecuada para entrenar o evaluar modelos de segmentaci\'on sobre resonancias magn\'eticas lumbares sagitales, incluyendo estructuras alineadas con el alcance del prototipo: v\'ertebras, discos intervertebrales y canal espinal. Esto permite construir una prueba de concepto t\'ecnicamente viable y reproducible.

Adem\'as, existen arquitecturas y frameworks de aprendizaje profundo ampliamente utilizados en segmentaci\'on m\'edica, como U-Net, nnU-Net y herramientas especializadas en im\'agenes biom\'edicas. Estos antecedentes permiten enfocar el valor del proyecto no en crear necesariamente una arquitectura completamente nueva, sino en integrar un flujo completo: carga de estudios, preprocesamiento, segmentaci\'on, medici\'on, visualizaci\'on y salida editable.

Desde el punto de vista acad\'emico, el proyecto resulta pertinente para un Proyecto Final de Ingenier\'ia porque combina investigaci\'on aplicada, dise\~no de software, procesamiento de datos, inteligencia artificial, validaci\'on t\'ecnica y documentaci\'on de resultados. La propuesta no se limita a entrenar un modelo aislado, sino que busca construir un prototipo funcional con componentes articulados y verificables. Esto permite aplicar conocimientos propios de la carrera en un problema concreto y con restricciones reales.

A su vez, el proyecto se plantea dentro de l\'imites claros. No se propone desarrollar un sistema de diagn\'ostico aut\'onomo, ni una herramienta certificada para uso cl\'inico, ni una plataforma hospitalaria completa. Esta delimitaci\'on es importante porque evita sobreprometer capacidades y permite concentrar el esfuerzo en un MVP t\'ecnicamente alcanzable: segmentar estructuras, obtener mediciones simples y presentar resultados revisables.

Desde el punto de vista funcional, la herramienta propuesta puede aportar valor como soporte al an\'alisis estructural de RM lumbar. La segmentaci\'on autom\'atica permitir\'ia identificar visualmente regiones de inter\'es; las mediciones geom\'etricas aportar\'ian informaci\'on cuantitativa derivada de las m\'ascaras; y la salida estructurada editable permitir\'ia organizar los resultados para revisi\'on profesional. En todos los casos, el profesional conserva el control final sobre la interpretaci\'on, correcci\'on y eventual descarte de los resultados.

El estado del arte muestra que existen avances relevantes en segmentaci\'on autom\'atica de columna lumbar, datasets p\'ublicos y soluciones comerciales orientadas a mediciones y reportes. Sin embargo, muchos antecedentes se concentran en componentes parciales: algunos aportan datos, otros modelos de segmentaci\'on, otros m\'etricas de evaluaci\'on y otros productos cerrados. La brecha que justifica este PFI se encuentra en integrar esos elementos en un prototipo acad\'emico, acotado, reproducible y orientado a revisi\'on profesional.

El proyecto tambi\'en se justifica por su enfoque de trazabilidad. Cada medici\'on generada debe estar asociada a una estructura segmentada visible, y cada resultado debe poder ser revisado por el usuario. Esta caracter\'istica permite diferenciar el prototipo de una \enquote{caja negra} diagn\'ostica y orientarlo hacia una herramienta asistiva, donde la inteligencia artificial act\'ua como apoyo operativo.

Finalmente, la elecci\'on de trabajar inicialmente con datasets p\'ublicos reduce riesgos vinculados con privacidad y datos sensibles. Si en etapas futuras se incorporaran im\'agenes cl\'inicas no p\'ublicas, deber\'ian contemplarse procedimientos de anonimizaci\'on, consentimiento y resguardo normativo. Para el alcance actual, el uso de datos p\'ublicos permite concentrar el desarrollo en la validaci\'on t\'ecnica del pipeline y en la construcci\'on del prototipo funcional.

En s\'intesis, el proyecto se justifica porque aborda una necesidad real de asistencia en el an\'alisis estructural de RM lumbar, utiliza recursos t\'ecnicos disponibles, mantiene un alcance acad\'emico responsable y propone una integraci\'on funcional que conecta segmentaci\'on, medici\'on, visualizaci\'on y revisi\'on profesional.
